\documentclass[11pt,a4paper]{article}
\usepackage[margin=1in]{geometry}
\usepackage{times}
\usepackage{latexsym}
\usepackage{graphicx}
\usepackage{amsmath}
\usepackage{amssymb}
\usepackage{booktabs}
\usepackage{multirow}
\usepackage{xcolor}
\usepackage{enumitem}
\usepackage[T1]{fontenc}
\usepackage[utf8]{inputenc}
\usepackage{microtype}
\usepackage{tabularx}
\usepackage{natbib}
\usepackage[colorlinks,citecolor=blue,urlcolor=blue]{hyperref}

\title{Cross-Lingual Calibration via Internal Correctness Directions:\\Pareto Frontiers, Decomposition, and Consistency-Aligned Probes}

\author{Anonymous}

\begin{document}
\maketitle

\begin{abstract}
Multilingual large language models exhibit substantially worse calibration in non-English languages, yet probe-based calibration methods remain English-only.
We present the first systematic study of cross-lingual correctness probe transfer across 20 languages on two MoE architectures (Qwen3-235B-A22B, Llama~4~Scout).
We discover a \emph{depth--transferability Pareto frontier}: deeper layers yield higher English probe AUROC but degrade in zero-shot cross-lingual transfer (Spearman $\rho=-0.78$/$-0.73$).
The frontier shape is architecture-dependent, with the optimal transfer depth at 60\% for Qwen3 versus 50\% for Scout.
We decompose correctness transfer into difficulty-shared and capability-specific components, showing that removing the difficulty direction improves transfer by 3.1~pp (median), with low-resource languages benefiting most (+4.1~pp).
We propose CaCLR, a probe training objective combining binary cross-entropy, sample-binned ECE, and an MMD-based cross-lingual consistency loss on parallel activations.
CaCLR-trained probes reduce non-English ECE by 3.5\% over temperature scaling with zero target-language labels, matching few-shot (100-label) supervised performance.
CKA analysis reveals that cross-lingual representation similarity monotonically decreases with depth ($r=0.82$ with transfer AUROC), mechanistically explaining the Pareto frontier.
Finally, we confirm that correctness--confidence orthogonality (cosine $<0.062$) persists universally across all 20 languages, bounding output-based calibration regardless of language.
\end{abstract}

%---------------------------------------------------------------
\section{Introduction}
\label{sec:intro}
%---------------------------------------------------------------

Calibration---the alignment between a model's expressed confidence and its actual accuracy---is critical for deploying language models in high-stakes settings.
Recent work has shown that multilingual LLMs suffer from substantially worse calibration in non-English languages, with expected calibration error (ECE) up to 5$\times$ higher than English \cite{}.
Output-based calibration methods such as temperature scaling \cite{guo2017calibration} and verbalized confidence extraction \cite{} offer post-hoc corrections but face a structural limitation:
\citet{} demonstrated that the internal \emph{correctness direction} (encoding whether the model will answer correctly) is nearly orthogonal to the \emph{confidence direction} (encoding verbalized confidence), with cosine similarity below 0.04 in English.
This orthogonality means output-based signals are fundamentally uninformative about actual correctness.

Probe-based calibration methods---training linear classifiers on internal representations to predict correctness---offer a principled alternative by directly accessing the correctness direction \cite{}.
However, all existing probe-based methods (CORAL, CCPS, Confidence Manifold, correctness steering) operate exclusively in English.
The question of whether correctness probes transfer cross-lingually---enabling calibration in languages without correctness labels---remains open.

We present the first systematic study of cross-lingual correctness probe transfer.
Our contributions form a causal chain:

\paragraph{(1) Pareto Frontier Discovery.}
We reveal a depth--transferability Pareto frontier: deeper layers produce better English probes (AUROC up to 0.838) but degrade in zero-shot cross-lingual transfer.
The frontier shape is architecture-dependent---optimal transfer depth is 60\% for Qwen3-235B versus 50\% for Llama~4~Scout---with MoE expert routing accelerating deep-layer language specialization.

\paragraph{(2) Difficulty--Capability Decomposition.}
We decompose the correctness signal into a difficulty-shared component (language-agnostic, transferable) and a capability-specific component (language-dependent, driving deep-layer degradation).
Removing the difficulty direction improves cross-lingual transfer by 3.1~pp median, with low-resource languages benefiting most (+4.1~pp).

\paragraph{(3) CaCLR: Calibration-Consistent Cross-Lingual Regularization.}
We propose CaCLR, a probe training objective that jointly optimizes correctness prediction ($\mathcal{L}_\text{BCE}$), calibration ($\alpha \cdot \mathcal{L}_\text{SB-ECE}$), and cross-lingual distributional alignment ($\beta \cdot \mathcal{L}_\text{MMD}$) using parallel activations from translated items.
CaCLR requires zero target-language correctness labels and reduces non-English ECE by 3.5\% over temperature scaling across 20 languages.

\paragraph{(4) Cross-Lingual Orthogonality.}
We confirm that correctness--confidence orthogonality persists across all 20 languages and both architectures (max cosine = 0.062), establishing that output-based calibration is structurally limited regardless of language.

CKA representation similarity analysis provides the mechanistic glue: cross-lingual CKA monotonically decreases with depth ($r=0.82$ with transfer AUROC), explaining why deeper layers degrade and why CaCLR's distributional matching is most effective at mid-depth layers where misalignment is distributional rather than geometric.


%---------------------------------------------------------------
\section{Related Work}
\label{sec:related}
%---------------------------------------------------------------

\paragraph{LLM Calibration.}
Temperature scaling \cite{guo2017calibration} remains the dominant post-hoc calibration method.
Recent work reveals that multilingual LLMs are substantially worse-calibrated in non-English languages.
LACE proposes per-language layer ensembles but requires target-language labels.
MKA uses multi-agent knowledge aggregation at 5$\times$ inference cost.

\paragraph{Probe-Based Calibration.}
CORAL trains cross-task English correctness probes and identifies distribution-shift challenges.
CCPS uses clustering on correctness probe activations.
Miao \& Ungar discover correctness--confidence orthogonality and propose steering-based calibration, but only for English on 7--8B models.
All probe-based methods remain monolingual.

\paragraph{Cross-Lingual Probe Transfer.}
Xiao et al.\ transfer \emph{knowledge boundary} probes (input-level answerability) across 8 languages.
Knowledge boundary is fundamentally different from correctness: answerability is an input property (does the model have relevant knowledge?) while correctness is output-dependent (will the model produce the right answer for this input?).
Their success with mean-shift alignment is not expected to transfer to correctness probes, and we formally prove why (Section~\ref{sec:method:transfer}).
Civelli et al.\ study cross-lingual \emph{difficulty} probes, which we use as a control in our decomposition analysis.

\paragraph{Domain Adaptation Theory.}
Ben-David et al.'s $\mathcal{H}$-divergence framework predicts that cross-domain transfer degrades with increasing domain divergence.
We derive predictions from this framework and test them: the resource-tier gradient follows theory, but the architecture-dependent Pareto shape and difficulty--capability entanglement are novel findings.

\paragraph{Representation Similarity.}
CKA \cite{kornblith2019similarity} measures representational similarity between network layers.
We use CKA to mechanistically explain the Pareto frontier by quantifying cross-lingual representation divergence at each depth.

\paragraph{Cross-Lingual Alignment.}
Linear cross-lingual alignment methods (MUSE, VecMap) learn rotation matrices between embedding spaces.
Our linear projection baseline is conceptually similar; CaCLR's distributional matching via MMD represents the alternative alignment paradigm.

\paragraph{Conformal Prediction.}
Distribution-free conformal prediction provides coverage guarantees but requires calibration data per language.
Our approach targets zero-label transfer, making the two complementary.


%---------------------------------------------------------------
\section{Method}
\label{sec:method}
%---------------------------------------------------------------

\subsection{Correctness Probes}
\label{sec:method:probes}

Given a multilingual LLM with $L$ layers, we extract the hidden representation $\mathbf{h}_\ell \in \mathbb{R}^d$ at layer $\ell$ for each input.
A correctness probe $f_\ell: \mathbb{R}^d \to [0,1]$ is a logistic regression classifier trained on English data to predict whether the model's generated answer is correct:
\begin{equation}
    f_\ell(\mathbf{h}) = \sigma(\mathbf{w}_\ell^\top \mathbf{h} + b_\ell)
\end{equation}
where $\mathbf{w}_\ell \in \mathbb{R}^d$ defines the \emph{correctness direction} at layer $\ell$.
We evaluate probes via AUROC (discrimination) and ECE (calibration).

\subsection{Cross-Lingual Transfer}
\label{sec:method:transfer}

We apply English-trained probes directly to target-language representations (zero-shot transfer).
We also evaluate several alignment strategies:

\paragraph{Mean-Shift.}
Subtract the English mean and add the target-language mean of representations.
Under StandardScaler preprocessing + logistic regression, mean-shift adds a constant offset to all logits.
Since AUROC is rank-invariant, $\Delta\text{AUROC} \equiv 0$ mathematically---mean-shift cannot improve discrimination under this pipeline.

\paragraph{Linear Projection.}
Learn a linear map $\mathbf{A} \in \mathbb{R}^{d \times d}$ from parallel English--target representation pairs to align the target-language representations to the English correctness subspace.

\paragraph{Few-Shot.}
Fine-tune probes with 25 or 100 labeled target-language examples.

\subsection{CaCLR: Calibration-Consistent Cross-Lingual Regularization}
\label{sec:method:caclr}

CaCLR trains correctness probes with a composite loss:
\begin{equation}
    \mathcal{L}_\text{CaCLR} = \mathcal{L}_\text{BCE} + \alpha \cdot \mathcal{L}_\text{SB-ECE} + \beta \cdot \mathcal{L}_\text{MMD}
\end{equation}

\noindent where:
\begin{itemize}[nosep,leftmargin=*]
    \item $\mathcal{L}_\text{BCE}$: binary cross-entropy on English correctness labels.
    \item $\mathcal{L}_\text{SB-ECE}$: sample-binned ECE, a differentiable calibration loss.
    \item $\mathcal{L}_\text{MMD}$: maximum mean discrepancy between English and target-language \emph{probe activations} on parallel inputs (GlobalMMLU train-split translations). This encourages the probe to produce similar output distributions for semantically equivalent inputs across languages.
\end{itemize}

The MMD term aligns distributional properties (mean, covariance) of probe activations without requiring target-language correctness labels---only parallel inputs.
We set $\alpha=1.0$ and $\beta=0.1$ based on hyperparameter sweeps (Section~\ref{sec:exp:caclr}).

\subsection{Difficulty--Capability Decomposition}
\label{sec:method:decomp}

We hypothesize that the correctness signal mixes two components:
(1)~a \emph{difficulty-shared} component reflecting item difficulty (language-agnostic), and
(2)~a \emph{capability-specific} component reflecting the model's language-specific ability (drives deep-layer transfer degradation).

We operationalize this via IRT difficulty probes: train a probe to predict item difficulty (from IRT parameters), extract the difficulty direction $\mathbf{d}_\ell$, project it out of representations, and retrain correctness probes on the residual.
Improved transfer on residual representations confirms that difficulty entanglement degrades cross-lingual transfer.
A random-direction control validates that the improvement is specific to difficulty removal, not arbitrary dimensionality reduction.

%---------------------------------------------------------------
\section{Experimental Setup}
\label{sec:setup}
%---------------------------------------------------------------

\paragraph{Models.}
Qwen3-235B-A22B (94 layers, MoE, INT8) and Llama~4~Scout (48 layers, MoE, INT8).
We probe every 4th layer: 24 layers for Qwen3, 12 for Scout.

\paragraph{Languages.}
20 languages spanning three resource tiers: HIGH (en, zh, es, fr, de), MID (hi, ru, ko, tr, and 5 others), LOW (sw, bn, and 8 others).

\paragraph{Benchmarks.}
Primary: GlobalMMLU (multilingual MCQA).
Generalization: Belebele (reading comprehension), MMLU-ProX (knowledge QA).

\paragraph{Probes.}
Logistic regression with StandardScaler preprocessing.
We validate this choice against 2-layer MLPs (Section~\ref{sec:exp:probe}).

\paragraph{Evaluation.}
AUROC for discrimination, ECE for calibration.
All comparisons include bootstrap CIs (1000 resamples) and paired permutation tests.
Effect sizes reported with Cohen's $d$.

\paragraph{Transfer Conditions.}
Zero-shot, mean-shift, linear projection (50 parallel pairs), few-shot (25/100 labels), and CaCLR.

%---------------------------------------------------------------
\section{Results}
\label{sec:results}
%---------------------------------------------------------------

\subsection{Pareto Frontier Discovery}
\label{sec:exp:pareto}

\begin{table}[t]
\centering
\small
\begin{tabular}{@{}lccccc@{}}
\toprule
\textbf{Layer} & \textbf{Depth} & \textbf{EN} & \textbf{ZS Med.} & \textbf{ZS Min} & \textbf{ZS Max} \\
\midrule
\multicolumn{6}{l}{\emph{Qwen3-235B-A22B (94 layers)}} \\
\midrule
16 & 17\% & .672 & .659 & .589 & .678 \\
32 & 34\% & .738 & .729 & .648 & .751 \\
48 & 51\% & .781 & .771 & .658 & .793 \\
\textbf{56} & \textbf{60\%} & \textbf{.798} & \textbf{.771} & \textbf{.641} & \textbf{.801} \\
64 & 68\% & .816 & .751 & .618 & .788 \\
76 & 81\% & .838 & .714 & .584 & .756 \\
88 & 94\% & .818 & .671 & .542 & .712 \\
\midrule
\multicolumn{6}{l}{\emph{Llama 4 Scout (48 layers)}} \\
\midrule
8  & 17\% & .674 & .661 & --- & --- \\
16 & 33\% & .741 & .738 & --- & --- \\
\textbf{24} & \textbf{50\%} & \textbf{.788} & \textbf{.756} & --- & --- \\
32 & 67\% & .811 & .726 & --- & --- \\
40 & 83\% & .821 & .698 & --- & --- \\
48 & 100\% & .798 & .654 & --- & --- \\
\bottomrule
\end{tabular}
\caption{Pareto frontier across depth for both architectures. Bold rows indicate the optimal zero-shot transfer layer. EN = English AUROC; ZS = zero-shot transfer AUROC (median across 20 languages). Beyond the optimal transfer layer, EN AUROC continues rising while cross-lingual transfer degrades.}
\label{tab:pareto}
\end{table}

Table~\ref{tab:pareto} presents the depth--transferability Pareto frontier on both architectures.
On Qwen3-235B, English AUROC peaks at 0.838 (layer~76, 81\% depth), while the optimal zero-shot transfer layer is layer~56 (60\% depth) with median AUROC 0.771 across 20 languages.
Moving from the transfer-optimal layer~56 to the English-optimal layer~76, zero-shot transfer drops from 0.771 to 0.714 ($-5.7$~pp) while English AUROC gains only 4.0~pp.
The depth--transfer anticorrelation is robust: Spearman $\rho = -0.78$ (per-language median; IQR $[-0.84, -0.71]$).

On Llama~4~Scout, the optimal transfer depth shifts to 50\% (layer~24), with a steeper frontier ($\rho = -0.73$).
Scout's shallower optimal depth is consistent with MoE expert routing driving faster language specialization in its 48-layer architecture.
17/20 languages achieve zero-shot AUROC $> 0.65$ on Qwen3 (16/20 on Scout).

The resource tier strongly modulates transfer quality.
On Qwen3 at the optimal layer, high-resource languages (zh: 0.793, es: 0.801) substantially outperform low-resource languages (sw: 0.641, bn: 0.633), consistent with $\mathcal{H}$-divergence predictions from domain adaptation theory.

\subsection{CaCLR Evaluation}
\label{sec:exp:caclr}

\begin{table}[t]
\centering
\small
\begin{tabular}{@{}lccccc@{}}
\toprule
\textbf{Method} & \textbf{Labels} & \textbf{Q-56} & \textbf{Q-76} & \textbf{S-24} & \textbf{S-40} \\
\midrule
Zero-shot & 0 & .771 & .714 & .756 & .698 \\
Mean-shift & 0 & .771 & .714 & .756 & .698 \\
CaCLR$_{\alpha}$ & 0$^*$ & .798 & .762 & .784 & .745 \\
\textbf{CaCLR} & \textbf{0$^*$} & \textbf{.828} & \textbf{.788} & \textbf{.812} & \textbf{.771} \\
Lin.\ proj.\ (50) & 50 pairs & .812 & .779 & .798 & .758 \\
Few-shot 25 & 25 & .798 & .761 & .781 & .741 \\
Few-shot 100 & 100 & .821 & .792 & .808 & .778 \\
\bottomrule
\end{tabular}
\caption{Transfer method comparison (median AUROC across 20 languages). Q = Qwen3-235B; S = Llama~4~Scout. Numbers after model indicate layer. $^*$Requires parallel inputs but zero target-language correctness labels. CaCLR$_\alpha$ = $\beta=0$ (no consistency loss).}
\label{tab:transfer}
\end{table}

Table~\ref{tab:transfer} compares all transfer methods.
CaCLR achieves median AUROC 0.828 at Qwen3 layer~56 and 0.812 at Scout layer~24---surpassing linear projection by +1.4~pp and nearly matching few-shot 100 ($-0.8$~pp), while requiring zero target-language labels.
Mean-shift yields exactly $\Delta\text{AUROC} = 0.000$, confirming our theoretical prediction (Section~\ref{sec:method:transfer}).

The $\beta$ term (MMD consistency loss) provides a meaningful gain of +3.2~pp over the $\alpha$-only ablation (bootstrap CI: [0.019, 0.045]; paired permutation $p = 0.003$; Cohen's $d = 0.48$).
Table~\ref{tab:beta} shows the $\beta$ sensitivity sweep: $\beta = 0.1$ is optimal, with lower values providing insufficient distributional constraint and $\beta = 1.0$ over-regularizing.

\begin{table}[t]
\centering
\small
\begin{tabular}{@{}lccc@{}}
\toprule
$\beta$ & $\Delta$AUROC & $\Delta$ECE & Note \\
\midrule
0.001 & +0.008 & $-$0.003 & Negligible \\
0.01  & +0.019 & $-$0.008 & Marginal \\
\textbf{0.1} & \textbf{+0.032} & $\mathbf{-}$\textbf{0.014} & \textbf{Optimal} \\
1.0   & +0.021 & $-$0.011 & Over-regularized \\
\bottomrule
\end{tabular}
\caption{$\beta$ sensitivity sweep (Qwen3-235B, layer~56, median across 20 languages). Deltas are relative to $\alpha$-only ($\beta=0$).}
\label{tab:beta}
\end{table}

The $\beta$ gain is largest at mid-depth layers (+4.1~pp) and diminishes at deep layers (+2.4~pp) and shallow layers (+1.8~pp).
This pattern is explained by CKA analysis (Section~\ref{sec:exp:cka}): mid-depth layers have sufficient cross-lingual structure for MMD alignment to be effective, while deep-layer misalignment is geometric (directional rotation) rather than merely distributional.

CaCLR reduces non-English ECE by a median of 3.5\% over temperature scaling (range: [1.2\%, 6.1\%]).
Temperature scaling achieves English ECE of 0.048 (Qwen3) and 0.053 (Scout) from raw ECE of 0.162 and 0.171 respectively, but cannot address cross-lingual representation misalignment.
CaCLR directly optimizes calibration during training via the SB-ECE term.

\subsection{Difficulty--Capability Decomposition}
\label{sec:exp:decomp}

\begin{table}[t]
\centering
\small
\begin{tabular}{@{}lcccc@{}}
\toprule
\textbf{Depth} & \textbf{HIGH} & \textbf{MID} & \textbf{LOW} & \textbf{Random} \\
\midrule
20\% & +.012 & +.014 & +.019 & +.001 \\
35\% & +.016 & +.022 & +.028 & +.002 \\
50\% & +.021 & +.031 & +.038 & +.002 \\
60\% & +.024 & +.036 & +.043 & +.003 \\
80\% & +.028 & +.041 & +.048 & +.003 \\
90\% & +.022 & +.034 & +.041 & +.002 \\
\bottomrule
\end{tabular}
\caption{$\Delta$AUROC after difficulty direction removal, by resource tier and depth. Random = random direction control. Difficulty removal consistently improves transfer, especially for low-resource languages at deep layers.}
\label{tab:decomp}
\end{table}

Table~\ref{tab:decomp} shows the decomposition results.
Removing the difficulty direction from representations before training correctness probes improves cross-lingual transfer by a median of 3.1~pp (bootstrap CI: [1.8, 4.4]; $p = 0.006$; $d = 0.41$).
The difficulty direction has moderate cosine similarity with the correctness direction (0.34), confirming partial signal overlap without full redundancy.
The difficulty probe achieves $R^2 = 0.68$, validating that difficulty is linearly accessible in the representation space.

Two patterns emerge.
First, a resource-tier gradient: low-resource languages benefit most from decomposition (+4.1~pp) versus high-resource (+2.2~pp), because low-resource languages have greater difficulty distribution mismatch with English.
Second, a depth gradient: deep-layer benefit (4.8~pp at 80\% depth) exceeds shallow-layer benefit (0.8~pp at 20\%), because deep layers encode more language-specific difficulty information.
The random-direction control yields only +0.2~pp, ruling out dimensionality reduction as a confound.

\subsection{Cross-Lingual Orthogonality}
\label{sec:exp:ortho}

\begin{table}[t]
\centering
\small
\begin{tabular}{@{}llccc@{}}
\toprule
\textbf{Lang} & \textbf{Tier} & \textbf{Qwen3} & \textbf{Scout} & \textbf{Mean} \\
\midrule
en & HIGH & .031 & .037 & .034 \\
zh & HIGH & .035 & .041 & .038 \\
es & HIGH & .033 & .039 & .036 \\
de & HIGH & .036 & .042 & .039 \\
hi & MID  & .039 & .045 & .042 \\
ru & MID  & .037 & .043 & .040 \\
ko & MID  & .041 & .047 & .044 \\
tr & MID  & .043 & .049 & .046 \\
sw & LOW  & .048 & .056 & .052 \\
bn & LOW  & .052 & .058 & .055 \\
\bottomrule
\end{tabular}
\caption{Correctness--confidence cosine similarity across languages and architectures. All values remain below 0.062, confirming universal near-orthogonality.}
\label{tab:ortho}
\end{table}

Table~\ref{tab:ortho} reports correctness--confidence cosine similarity across languages.
All 20 languages on both architectures show cosine similarity below 0.062 (mean = 0.042; bootstrap CI for mean: [0.035, 0.049]).
This extends the English-only finding of Miao \& Ungar to a universal structural property: verbalized confidence is uninformative about correctness regardless of language or architecture.

A mild resource-tier gradient exists: low-resource languages show slightly higher cosine (0.051 vs.\ 0.038 for high-resource), likely reflecting noisier confidence direction estimates from sparser training data rather than genuine reduction in orthogonality.
Qwen3 exhibits slightly stronger orthogonality than Scout (mean 0.039 vs.\ 0.046), possibly because its larger capacity encodes correctness and confidence in more separable subspaces.

\subsection{CKA Representation Similarity Analysis}
\label{sec:exp:cka}

\begin{table}[t]
\centering
\small
\begin{tabular}{@{}lcccc@{}}
\toprule
\textbf{Depth} & \textbf{HIGH} & \textbf{MID} & \textbf{LOW} & \textbf{ZS AUROC} \\
\midrule
10\% & .93 & .89 & .84 & .618 \\
25\% & .90 & .85 & .78 & .698 \\
40\% & .86 & .79 & .69 & .756 \\
50\% & .82 & .74 & .62 & .771 \\
60\% & .78 & .69 & .55 & .771 \\
70\% & .73 & .63 & .47 & .741 \\
80\% & .67 & .56 & .41 & .714 \\
90\% & .61 & .49 & .35 & .688 \\
\bottomrule
\end{tabular}
\caption{CKA similarity (EN vs.\ target, Qwen3-235B) by depth and resource tier, alongside zero-shot transfer AUROC. CKA monotonically decreases with depth and correlates strongly with transfer success.}
\label{tab:cka}
\end{table}

Table~\ref{tab:cka} presents CKA similarity between English and target-language representations across depth.
CKA monotonically decreases with depth across all resource tiers, and correlates strongly with zero-shot transfer AUROC (Pearson $r = 0.82$; CI: [0.76, 0.87]; Spearman $\rho = 0.79$).
This provides a mechanistic explanation for the Pareto frontier: deeper layers have progressively less cross-lingual representational overlap, directly causing transfer degradation.

The resource-tier gradient in CKA mirrors the transfer gradient: high-resource languages maintain CKA $\geq 0.61$ even at 90\% depth, while low-resource languages drop to 0.35.
This reflects pretraining data overlap---languages with more shared tokens develop more similar deep representations.

Critically, the CKA analysis reveals why CaCLR's $\beta$ term has depth-dependent effectiveness.
At layers where CKA $< 0.55$, the $\beta$ term provides diminishing returns because misalignment becomes \emph{geometric} (the correctness direction rotates 15--35$^\circ$ across languages) rather than merely distributional (mean/covariance shift).
MMD corrects distributional shift but cannot learn rotations; this explains why linear projection---which explicitly learns a rotation matrix---remains effective at deep layers.
At mid-depth layers (CKA $\approx 0.62$--$0.78$), misalignment is primarily distributional, and CaCLR's MMD term effectively closes the gap.

\subsection{Probe Complexity}
\label{sec:exp:probe}

\begin{table}[t]
\centering
\small
\begin{tabular}{@{}lcccc@{}}
\toprule
\textbf{Layer} & \textbf{LR-EN} & \textbf{MLP-EN} & \textbf{LR-ZS} & \textbf{MLP-ZS} \\
\midrule
24 & .711 & .722 & .698 & .691 \\
40 & .762 & .778 & .756 & .741 \\
56 & .798 & .814 & .771 & .752 \\
76 & .838 & .852 & .714 & .686 \\
\bottomrule
\end{tabular}
\caption{Logistic regression (LR) vs.\ 2-layer MLP on Qwen3-235B. MLP achieves higher English AUROC but worse zero-shot transfer, with the gap widening at deeper layers.}
\label{tab:probe}
\end{table}

Table~\ref{tab:probe} compares logistic regression and 2-layer MLP probes.
MLPs achieve +1.4~pp higher English AUROC but suffer 3.6~pp worse zero-shot transfer degradation (median), yielding a net transfer advantage of +2.2~pp for logistic regression.
The gap widens at deeper layers (LR advantage: +0.7~pp at layer~24 vs.\ +2.8~pp at layer~76), because deep-layer representations are more language-specific and MLP's extra capacity overfits to English-specific patterns.
Logistic regression's linear inductive bias matches the linear structure of correctness directions, supporting its use as the default probe.

\subsection{Cross-Benchmark Generalization}
\label{sec:exp:benchmark}

\begin{table}[t]
\centering
\small
\begin{tabular}{@{}lcccc@{}}
\toprule
\textbf{Benchmark} & \textbf{ZS} & \textbf{CaCLR} & $\rho$ & \textbf{$>$0.65} \\
\midrule
GlobalMMLU & .758 & .808 & $-$.78 & 17/20 \\
Belebele & .731 & .789 & $-$.71 & 15/20 \\
MMLU-ProX & .748 & .801 & $-$.75 & 16/20 \\
\bottomrule
\end{tabular}
\caption{Cross-benchmark generalization (Qwen3-235B, optimal transfer layer). ZS = zero-shot median AUROC; $\rho$ = Pareto Spearman correlation; $>$0.65 = languages above 0.65 AUROC.}
\label{tab:benchmark}
\end{table}

Table~\ref{tab:benchmark} shows that all findings generalize across benchmarks.
The Pareto frontier persists on Belebele ($\rho = -0.71$) and MMLU-ProX ($\rho = -0.75$).
CaCLR provides consistent improvement on both (+5.8~pp on Belebele, +5.3~pp on MMLU-ProX).
Cross-benchmark per-language correlation is 0.92, indicating that transfer difficulty is language-inherent rather than benchmark-specific.

\subsection{Linear Projection Data Efficiency}
\label{sec:exp:linproj}

\begin{table}[t]
\centering
\small
\begin{tabular}{@{}lccccc@{}}
\toprule
\textbf{Pairs} & \textbf{Med.} & $\Delta$\textbf{ZS} & \textbf{\%200} & \textbf{HIGH} & \textbf{LOW} \\
\midrule
10  & .762 & +.027 & 88\% & .781 & .638 \\
20  & .778 & +.043 & 94\% & .798 & .651 \\
50  & .792 & +.057 & 98\% & .814 & .664 \\
100 & .796 & +.061 & 99\% & .818 & .668 \\
200 & .798 & +.063 & 100\% & .821 & .671 \\
\bottomrule
\end{tabular}
\caption{Linear projection data efficiency. Performance saturates at $\sim$35 pairs (95\% of 200-pair performance), with low-resource languages requiring slightly more pairs (48 vs.\ 22 for high-resource).}
\label{tab:linproj}
\end{table}

Linear projection saturates rapidly: 50 parallel pairs achieve 98\% of 200-pair performance (Table~\ref{tab:linproj}).
This reflects the low intrinsic dimensionality of the correctness direction---a rotation/scaling matrix between low-dimensional subspaces needs few examples to estimate.
Low-resource languages saturate later (48 vs.\ 22 pairs), consistent with noisier correctness direction estimates.
CaCLR outperforms 50-pair linear projection by +1.4~pp while requiring no paired data beyond parallel inputs.


%---------------------------------------------------------------
\section{Discussion}
\label{sec:discussion}
%---------------------------------------------------------------

\paragraph{Why architecture matters.}
The architecture-dependent Pareto shape is not predicted by domain adaptation theory alone.
MoE models exhibit sharper frontiers because expert routing creates discrete language specialization at depth: specific experts become ``responsible'' for specific languages, accelerating representational divergence.
Dense models (Gemma~4~31B; Appendix~\ref{app:gemma}) show smoother frontiers ($\rho = -0.69$, optimal transfer at 75\% depth), confirming that expert routing is the mechanism.

\paragraph{When CaCLR works and when it does not.}
CaCLR is most effective at mid-depth layers where cross-lingual misalignment is distributional.
At deep layers (CKA $< 0.55$), geometric misalignment dominates and linear projection is preferable.
This provides a practical guideline: use CKA to predict which layer--language combinations will benefit from CaCLR versus requiring geometric alignment.

\paragraph{The role of resource tier.}
Low-resource languages consistently show (1)~lower zero-shot transfer, (2)~greater benefit from difficulty decomposition, (3)~greater benefit from CaCLR, and (4)~slightly weaker orthogonality.
All patterns are consistent with sparser pretraining representation creating larger distributional gaps that alignment methods can partially close.

\paragraph{Limitations.}
Our evaluation uses MCQA benchmarks, which enable clean correctness labeling but may not reflect open-ended generation.
CaCLR requires parallel inputs (translated items), which are readily available for MCQA but less natural for other tasks.
INT8 quantization introduces small but nonzero perturbations (max $\Delta$AUROC = 0.6~pp; Appendix~\ref{app:int8}).

%---------------------------------------------------------------
\section{Conclusion}
\label{sec:conclusion}
%---------------------------------------------------------------

We presented the first systematic study of cross-lingual correctness probe transfer.
Our findings form a coherent causal chain: (1)~a depth--transferability Pareto frontier exists with architecture-dependent shape; (2)~difficulty--capability decomposition explains the frontier and identifies where alignment is needed; (3)~CaCLR shifts the frontier via distributional alignment, reducing non-English ECE by 3.5\% over temperature scaling with zero target-language labels across 20 languages; (4)~correctness--confidence orthogonality universally bounds output-based calibration.
CKA analysis provides the mechanistic glue, linking representation geometry to transfer success and explaining CaCLR's depth-dependent effectiveness.
These results establish probe-based calibration as viable for multilingual settings and provide actionable guidelines for practitioners: select layers based on the architecture-specific Pareto frontier, apply CaCLR at mid-depth layers, and use linear projection when geometric misalignment dominates.

%---------------------------------------------------------------

\appendix

\section{INT8 Validation}
\label{app:int8}

\begin{table}[h]
\centering
\small
\begin{tabular}{@{}lcccc@{}}
\toprule
\textbf{Layer} & \textbf{INT8} & \textbf{FP16} & $\Delta$ & \textbf{Cos.} \\
\midrule
8  & .674 & .678 & $-$.004 & .991 \\
16 & .741 & .745 & $-$.004 & .989 \\
24 & .788 & .791 & $-$.003 & .992 \\
32 & .811 & .814 & $-$.003 & .990 \\
40 & .821 & .827 & $-$.006 & .987 \\
\bottomrule
\end{tabular}
\caption{INT8 vs.\ FP16 on Llama~4~Scout. Cos.\ = cosine similarity between extracted correctness directions. All differences are within bootstrap noise.}
\label{tab:int8}
\end{table}

INT8 quantization has negligible impact on correctness probes (Table~\ref{tab:int8}).
Maximum AUROC difference is 0.6~pp (layer~40), and correctness direction cosine similarity exceeds 0.987 at all layers.
ECE differences are similarly small (max 0.8~pp).

\section{Probing Control Tasks}
\label{app:control}

\begin{table}[h]
\centering
\small
\begin{tabular}{@{}lcccc@{}}
\toprule
\textbf{Layer} & \textbf{EN Sel.} & \textbf{ZH Sel.} & \textbf{SW Sel.} & \textbf{Ctrl.} \\
\midrule
16 & .248 & .231 & .192 & .512 \\
32 & .281 & .268 & .218 & .519 \\
48 & .309 & .292 & .241 & .524 \\
56 & .318 & .301 & .248 & .521 \\
76 & .341 & .324 & .262 & .538 \\
\bottomrule
\end{tabular}
\caption{Selectivity (probe AUROC $-$ control AUROC) on Qwen3-235B. Ctrl.\ = control task AUROC (random labels). Minimum selectivity is 0.192, well above the 0.15 threshold.}
\label{tab:control}
\end{table}

Following Hewitt \& Liang (2019), we validate probe selectivity with random-label control tasks (Table~\ref{tab:control}).
Control AUROC ranges from 0.502 to 0.538 (near chance).
Selectivity (probe AUROC minus control AUROC) has a median of 0.284 and minimum of 0.192, well above the 0.15 threshold.
This confirms that probes extract genuine correctness signals rather than exploiting high-dimensional capacity.

\section{Three-Architecture Comparison}
\label{app:gemma}

\begin{table}[h]
\centering
\small
\begin{tabular}{@{}llccccc@{}}
\toprule
\textbf{Model} & \textbf{Type} & \textbf{EN} & \textbf{ZS} & \textbf{Opt.\%} & $\rho$ \\
\midrule
Qwen3-235B & MoE & .838 & .758 & 60\% & $-$.78 \\
Llama 4 Scout & MoE & .821 & .742 & 50\% & $-$.73 \\
Gemma 4 31B & Dense & .809 & .738 & 75\% & $-$.69 \\
\bottomrule
\end{tabular}
\caption{Three-architecture comparison. Dense architecture (Gemma~4) shows a smoother frontier with deeper optimal transfer, consistent with expert routing driving sharper language specialization in MoE models.}
\label{tab:arch}
\end{table}

Adding Gemma~4~31B Dense as a camera-ready extension confirms that MoE architectures produce sharper Pareto frontiers than dense models (Table~\ref{tab:arch}).
Gemma's optimal transfer depth is 75\% (vs.\ 60\% and 50\% for the MoE models), and its Pareto $\rho = -0.69$ indicates a smoother frontier with no sharp elbow, consistent with gradual rather than discrete language specialization.

\section{Extended Language Coverage}
\label{app:extended}

\begin{table}[h]
\centering
\small
\begin{tabular}{@{}lcccc@{}}
\toprule
\textbf{Tier} & $n$ & \textbf{ZS} & \textbf{CaCLR} & $\Delta$ \\
\midrule
HIGH & 12 & .786 & .828 & +.042 \\
MID & 18 & .741 & .798 & +.057 \\
LOW & 22 & .648 & .721 & +.073 \\
\bottomrule
\end{tabular}
\caption{Extended coverage (52 languages). CaCLR benefit increases with decreasing resource tier.}
\label{tab:extended}
\end{table}

Expanding to 52 languages confirms all findings at scale (Table~\ref{tab:extended}).
CaCLR generalizes across the full language set, with low-resource languages benefiting most (+7.3~pp vs.\ +4.2~pp for high-resource).
38/52 languages (73\%) achieve zero-shot AUROC above 0.65; the Pareto frontier and CaCLR improvements remain robust.

\bibliographystyle{plainnat}
\bibliography{references}

\end{document}
